前面几个单元处理的是一般矩阵，得到的结论也就一般：特征值可能是复数，特征向量可能几乎平行，换基矩阵可能接近奇异，甚至凑不齐一组特征基。可你在数据工作里遇到的矩阵大多不是一般矩阵。协方差、Gram 矩阵、核矩阵、图的 Laplacian、Hessian，它们全都满足同一条 \(A^{\trans}=A\)——因为它们记录的是成对关系，而``\(i\) 对 \(j\)''与``\(j\) 对 \(i\)''本来就不分先后。

这一条看似轻描淡写的假设，换来的东西多得不成比例。本单元要讲清楚它换来了什么，以及这些好处怎样一路传到 PCA、优化收敛速度和高斯采样上去。

\begin{center}
\begin{tikzpicture}[x=1cm,y=1cm,font=\footnotesize,>=stealth]
  \draw[black!45,fill=blue!6] (0,0) rectangle (3.3,1.4);
  \node[align=center] at (1.65,0.7) {\(A^{\trans}=A\)\\[2pt] 一条对称假设};
  \draw[->,thick,black!60] (3.4,0.7) -- (3.85,0.7);
  \draw[black!45,fill=blue!6] (3.95,0) rectangle (7.25,1.4);
  \node[align=center] at (5.6,0.7) {\(A=Q\Lambda Q^{\trans}\)\\[2pt] 实特征值、正交主轴};
  \draw[->,thick,black!60] (7.35,0.7) -- (7.8,0.7);
  \draw[black!45,fill=blue!6] (7.9,0) rectangle (11.2,1.4);
  \node[align=center] at (9.55,0.7) {\(x^{\trans}Ax=\sum_i\lambda_iy_i^2\)\\[2pt] 碗、鞍或槽};
  \draw[->,thick,black!60] (11.3,0.7) -- (11.75,0.7);
  \draw[black!45,fill=blue!6] (11.85,0) rectangle (15.15,1.4);
  \node[align=center] at (13.5,0.7) {\(A=LL^{\trans}\)\\[2pt] 可判定、可解、可采样};
\end{tikzpicture}

\emph{本单元的主线只有一条链，三个箭头分别是谱定理、换到特征坐标、要求全部特征值为正。每一步都是等价改写而不是新增假设：对称性买来正交特征基，正交特征基把二次型化成一组独立的平方项，全部系数为正就是正定，而正定又允许把这组平方项用一次消元直接写出来。}
\end{center}

链条上的四个环节各对应一节。第一节证明谱定理并交代对称性用在了哪一步、拿掉它会怎样。第二节转到标量函数 \(x^{\trans}Ax\)，用特征值的符号给几何分类，并说明正定的几条判据为什么是同一句话；条件数在这里第一次露面，它把线性代数和优化的收敛速度接在一起。第三节的 Cholesky 分解既是正定性的构造性证明，也是从模型协方差生成相关高斯样本、计算 \(\log\det\) 与解正规方程的实用工具，顺带解释``矩阵不正定''这条报错在实践中意味着什么。第四节把前三节放回可见的几何：旋转拿走交叉项，平移找到中心，剩下的分类只看特征值符号。

\begin{itemize}
\tightlist
\item \hyperref[sec:03-Linear-Algebra/06-Symmetric-and-Positive-Definite-Matrices/01]{``对称矩阵与谱定理：什么时候特征方向一定正交''}：把``特征方向正交''从运气变成保证，并说清这条保证的代价是什么；
\item \hyperref[sec:03-Linear-Algebra/06-Symmetric-and-Positive-Definite-Matrices/02]{``二次型与正定性：每个方向上的能量是否为正''}：学会只看特征值符号就判断驻点身份、目标函数是否严格凸、参数是否唯一；
\item \hyperref[sec:03-Linear-Algebra/06-Symmetric-and-Positive-Definite-Matrices/03]{``Cholesky 分解：把正能量写成平方和''}：把正定性变成一次可执行的计算，并把分解失败读成一条诊断；
\item \hyperref[sec:03-Linear-Algebra/06-Symmetric-and-Positive-Definite-Matrices/04]{``二次曲面用谱分解显露几何形状''}：把误差椭圆、置信区域和优化等高线认成同一族曲面。
\end{itemize}

四节按依赖顺序排列，连着读最省力。若你是带着一个具体问题进来的——训练卡住了、协方差矩阵报错了、置信区域该怎么画——也可以直接打开最接近的那一节，再沿篇内链接回补前面的结论。需要的前置只有特征值与对角化，见 \hyperref[sec:03-Linear-Algebra/05-Eigenvalues-and-Eigenvectors/02]{``对角化与动力系统：把耦合演化变成独立缩放''}。
